\documentclass[a4paper,12pt]{article}

\usepackage[T2A]{fontenc}
\usepackage[utf8]{inputenc}
\usepackage[russian]{babel}

\usepackage{amsmath,amssymb,mathtools}
\usepackage{helvet}
\renewcommand{\familydefault}{\sfdefault}
\usepackage{icomma}
\usepackage{geometry}
\usepackage{microtype}
\usepackage{tikz}
\usetikzlibrary{arrows.meta,calc}
\usepackage{tabularx,booktabs}
\usepackage{enumitem}
\usepackage{hyperref}
\usepackage{parskip}
\usepackage{fancyhdr}
\usepackage{setspace}
\usepackage{xcolor}

\geometry{left=18mm,right=18mm,top=18mm,bottom=20mm}
\onehalfspacing

\definecolor{accent}{HTML}{008080}
\definecolor{darkaccent}{HTML}{004D4D}
\definecolor{textgray}{HTML}{333333}
\definecolor{softgray}{HTML}{F2F5F5}

\hypersetup{hidelinks}

\pagestyle{fancy}
\fancyhf{}
\fancyhead[L]{\textcolor{darkaccent}{Прикладные задачи}}
\fancyhead[R]{\textcolor{textgray}{ЕГЭ}}
\fancyfoot[C]{\thepage}

\setlist[itemize]{leftmargin=*,itemsep=0.25em,topsep=0.25em}
\setlist[enumerate]{leftmargin=*,itemsep=0.35em,topsep=0.35em}

\newcommand{\R}{\mathbb{R}}
\newcommand{\topic}{Прикладные задачи ЕГЭ: формулы, величины и отбор ответа}
\newcommand{\student}{Ксения Клыкова}
\newcommand{\teacher}{Лёвин Артём Александрович}

\begin{document}

\begin{titlepage}
\thispagestyle{empty}
\centering

\vspace*{35mm}

{\Large\textcolor{accent}{\textbf{Чек-лист}}\par}

\vspace{8mm}

{\Huge\textbf{\topic}\par}

\vspace{12mm}

{\large \teacher{} эксклюзивно для \student\par}

\vspace{8mm}

{\large \today\par}

\vfill

\begin{tikzpicture}[remember picture,overlay]
  \draw[
    line width=1.2pt,
    accent,
    opacity=0.65
  ]
  ([xshift=25mm,yshift=30mm]current page.south west)
  .. controls
  ([xshift=70mm,yshift=55mm]current page.south west)
  and
  ([xshift=-70mm,yshift=12mm]current page.south east)
  ..
  ([xshift=-25mm,yshift=38mm]current page.south east);
\end{tikzpicture}

\end{titlepage}

\section*{1. Главная идея прикладной задачи}

\textbf{Прикладная задача} — это задача, где реальная ситуация переводится на язык формул.

Обычно в условии уже дана формула. Нужно:

\begin{enumerate}
  \item внимательно прочитать текст;
  \item выписать данные;
  \item понять, какую величину нужно найти;
  \item привести единицы измерения к одному виду;
  \item подставить данные в формулу;
  \item аккуратно посчитать;
  \item выбрать именно тот ответ, который спрашивают.
\end{enumerate}

\textbf{Главный принцип:} прикладная задача чаще всего проверяет не только алгебру, но и внимательность к тексту.

\section*{2. Как оформлять «дано»}

В прикладных задачах полезно выписывать данные отдельно. Это делает решение более спокойным и надёжным.

Например:

\[
L_0=10\text{ м},
\qquad
\alpha=1{,}2\cdot10^{-5},
\qquad
\Delta L=4{,}5\text{ мм},
\qquad
T-?
\]

Если величины даны в разных единицах измерения, их нужно согласовать.

\[
4{,}5\text{ мм}=4{,}5\cdot10^{-3}\text{ м}.
\]

\textbf{Важно:} размерности записываются не всегда, но в задачах с физическим смыслом они часто помогают не ошибиться.

\section*{3. Что означает запись вида \(L(T)\)}

Запись

\[
L(T)
\]

означает, что величина \(L\) зависит от переменной \(T\).

Например:

\[
L(T)=L_0(1+\alpha T).
\]

Здесь:

\begin{itemize}
  \item \(L_0\) — постоянная величина;
  \item \(\alpha\) — постоянная величина;
  \item \(T\) — переменная;
  \item \(L(T)\) — значение длины при температуре \(T\).
\end{itemize}

В решении можно писать короче:

\[
L=L_0(1+\alpha T),
\]

если понятно, что \(L\) зависит от \(T\).

\section*{4. Изменение величины и символ \(\Delta\)}

Символ \(\Delta\) обозначает изменение величины.

Если рельс удлинился, то

\[
\Delta L=L-L_0.
\]

Отсюда:

\[
L=L_0+\Delta L.
\]

Если по формуле

\[
L=L_0(1+\alpha T),
\]

а с другой стороны

\[
L=L_0+\Delta L,
\]

то можно приравнять:

\[
L_0(1+\alpha T)=L_0+\Delta L.
\]

Раскрываем скобки:

\[
L_0+\alpha L_0T=L_0+\Delta L.
\]

Сокращаем \(L_0\):

\[
\alpha L_0T=\Delta L.
\]

Выражаем температуру:

\[
T=\frac{\Delta L}{\alpha L_0}.
\]

\section*{5. Работа со степенями десяти}

В прикладных задачах часто встречаются числа вида:

\[
1{,}2\cdot10^{-5},
\qquad
4{,}5\cdot10^{-3}.
\]

С десятичными дробями лучше работать аккуратно. Иногда удобно избавиться от запятой:

\[
4{,}5\cdot10^{-3}=45\cdot10^{-4}.
\]

Почему так:

\[
4{,}5=45\cdot10^{-1},
\]

поэтому

\[
4{,}5\cdot10^{-3}=45\cdot10^{-1}\cdot10^{-3}=45\cdot10^{-4}.
\]

\subsection*{Правило деления степеней}

\[
\frac{a^n}{a^k}=a^{n-k}.
\]

Например:

\[
\frac{10^{-4}}{10^{-5}}=10^{-4-(-5)}=10^1.
\]

\section*{6. Пример: тепловое расширение рельса}

При температуре \(0^\circ\text{C}\) рельс имеет длину \(L_0=10\) м. При нагревании его длина меняется по формуле:

\[
L=L_0(1+\alpha T),
\qquad
\alpha=1{,}2\cdot10^{-5}.
\]

На сколько градусов нужно нагреть рельс, чтобы он удлинился на \(4{,}5\) мм?

\subsection*{Решение}

Переведём миллиметры в метры:

\[
4{,}5\text{ мм}=4{,}5\cdot10^{-3}\text{ м}.
\]

Используем формулу:

\[
T=\frac{\Delta L}{\alpha L_0}.
\]

Подставляем данные:

\[
T=\frac{4{,}5\cdot10^{-3}}{10\cdot1{,}2\cdot10^{-5}}.
\]

Перепишем числитель без десятичной дроби:

\[
4{,}5\cdot10^{-3}=45\cdot10^{-4}.
\]

Знаменатель:

\[
10\cdot1{,}2\cdot10^{-5}=12\cdot10^{-5}.
\]

Тогда:

\[
T=\frac{45\cdot10^{-4}}{12\cdot10^{-5}}
=\frac{15}{4}\cdot10^{-4-(-5)}
=\frac{15}{4}\cdot10
=\frac{150}{4}
=37{,}5.
\]

Ответ:

\[
\boxed{37{,}5}
\]

\section*{7. Задачи с прибылью}

В экономических задачах часто нужно просто аккуратно подставить данные в формулу.

Например, месячная операционная прибыль задаётся формулой:

\[
P(Q)=Q(p-v)-F.
\]

Здесь:

\begin{itemize}
  \item \(Q\) — объём производства;
  \item \(p\) — цена одной единицы продукции;
  \item \(v\) — переменные затраты на одну единицу;
  \item \(F\) — постоянные расходы;
  \item \(P(Q)\) — прибыль при объёме \(Q\).
\end{itemize}

Если сказано, что прибыль должна быть \textbf{не меньше} некоторого значения, в типовых задачах на первичное освоение удобно приравнять прибыль к этому граничному значению.

\subsection*{Пример}

\[
p=600,
\qquad
v=300,
\qquad
F=700000,
\qquad
P=500000.
\]

\[
500000=Q(600-300)-700000.
\]

\[
500000=300Q-700000.
\]

\[
1200000=300Q.
\]

\[
Q=4000.
\]

Ответ:

\[
\boxed{4000}
\]

\section*{8. Задачи на изменение расстояния}

Иногда в условии не дана готовая формула для искомой величины. Тогда её нужно составить самостоятельно.

Например, глубина колодца вычисляется по формуле:

\[
H=5t^2.
\]

Если до дождя камень падал \(0{,}6\) с, а после дождя время уменьшилось на \(0{,}2\) с, то:

\[
t_1=0{,}6,
\qquad
t_2=0{,}6-0{,}2=0{,}4.
\]

Уровень воды поднялся на разность глубин:

\[
\Delta H=H_1-H_2.
\]

\[
\Delta H=5t_1^2-5t_2^2.
\]

Подставим:

\[
\Delta H=5\cdot0{,}6^2-5\cdot0{,}4^2.
\]

Удобно вынести общий множитель и применить формулу разности квадратов:

\[
\Delta H=5(0{,}6^2-0{,}4^2)
=5(0{,}6-0{,}4)(0{,}6+0{,}4).
\]

\[
\Delta H=5\cdot0{,}2\cdot1=1.
\]

Ответ:

\[
\boxed{1}
\]

\section*{9. Задачи с отбором корня}

Иногда после подстановки получается квадратное уравнение с двумя корнями. В ответ может идти только один из них.

\subsection*{Пример}

Высота полёта камня задаётся формулой:

\[
y=-\frac{1}{60}x^2+\frac76x.
\]

Стена имеет высоту \(9\) м. Камень должен пролететь на \(1\) м выше стены, значит:

\[
y=9+1=10.
\]

Подставим:

\[
10=-\frac{1}{60}x^2+\frac76x.
\]

Перенесём всё в одну сторону и умножим на \(60\):

\[
x^2-70x+600=0.
\]

Найдём корни:

\[
D=70^2-4\cdot600=4900-2400=2500.
\]

\[
x_1=\frac{70-50}{2}=10,
\qquad
x_2=\frac{70+50}{2}=60.
\]

Если спрашивают \textbf{наибольшее расстояние}, выбираем больший корень:

\[
\boxed{60}
\]

\section*{10. Типичные ошибки}

\begin{itemize}
  \item Не выписать данные и потерять важную величину из текста.
  \item Не перевести миллиметры в метры или наоборот.
  \item Перепутать, какая величина является переменной, а какая константой.
  \item Считать \(L(T)\) отдельной переменной вместо значения функции.
  \item Забыть, что \(\Delta L=L-L_0\).
  \item Ошибиться при работе со степенями десяти.
  \item Считать напрямую с десятичными дробями там, где проще вынести множитель или применить формулу сокращённого умножения.
  \item Не заметить, что данные могут быть заданы неявно.
  \item Получить два корня и записать оба, хотя в условии требуется наибольший или наименьший.
  \item Не проверить, какой ответ соответствует смыслу задачи.
\end{itemize}

\section*{11. Короткая памятка}

\begin{center}
\renewcommand{\arraystretch}{1.35}
\begin{tabularx}{\linewidth}{@{}p{0.32\linewidth}X@{}}
\toprule
\textbf{Шаг} & \textbf{Что сделать} \\
\midrule
1 & Прочитать текст и понять, какая величина требуется. \\
2 & Выписать данные в блок «дано». \\
3 & Привести единицы измерения к одному виду. \\
4 & Убрать лишнюю запись вида \(L(T)\), если она только показывает зависимость. \\
5 & Подставить данные в формулу. \\
6 & Аккуратно выразить неизвестную. \\
7 & Посчитать, используя степени десяти и сокращения. \\
8 & Если есть несколько корней, выбрать тот, который отвечает вопросу задачи. \\
\bottomrule
\end{tabularx}
\end{center}

\newpage

\section*{Тренировочный блок}

Решите задачи. В ответ запишите только число, если не указано иное.

\subsection*{Средний уровень}

\begin{enumerate}
  \item При температуре \(0^\circ\text{C}\) длина стержня равна \(20\) м. При нагревании его длина меняется по формуле
  \[
  L=L_0(1+\alpha T),
  \qquad
  \alpha=2\cdot10^{-5}.
  \]
  На сколько градусов нужно нагреть стержень, чтобы он удлинился на \(8\) мм?

  \item Прибыль предприятия задаётся формулой
  \[
  P(Q)=Q(p-v)-F.
  \]
  Найдите \(Q\), если
  \[
  P=300000,\quad p=500,\quad v=200,\quad F=600000.
  \]

  \item Глубина колодца вычисляется по формуле
  \[
  H=5t^2.
  \]
  До дождя камень падал \(0{,}8\) с, после дождя — \(0{,}6\) с. На сколько метров поднялся уровень воды?
\end{enumerate}

\subsection*{Выше среднего}

\begin{enumerate}[resume]
  \item Длина детали меняется по формуле
  \[
  L=L_0(1+\alpha T).
  \]
  Найдите \(T\), если
  \[
  L_0=5,\quad \alpha=4\cdot10^{-5},\quad \Delta L=3\text{ мм}.
  \]

  \item Прибыль фирмы задаётся формулой
  \[
  P(Q)=Q(p-v)-F.
  \]
  Найдите наименьший объём \(Q\), при котором прибыль равна \(800000\), если
  \[
  p=750,\quad v=350,\quad F=400000.
  \]

  \item Высота полёта тела задаётся формулой
  \[
  y=-\frac{1}{20}x^2+2x.
  \]
  На каком расстоянии \(x\) высота равна \(15\)? В ответ запишите меньшее значение \(x\).

  \item Высота полёта тела задаётся формулой
  \[
  y=-\frac{1}{40}x^2+\frac32x.
  \]
  Найдите наибольшее расстояние \(x\), при котором высота равна \(20\).

  \item Глубина вычисляется по формуле
  \[
  H=5t^2.
  \]
  Время падения уменьшилось с \(1{,}2\) с до \(0{,}9\) с. На сколько метров уменьшилась вычисляемая глубина?

  \item В формуле
  \[
  A=\frac{B}{CD}
  \]
  найдите \(B\), если
  \[
  A=2{,}5\cdot10^3,\quad C=5\cdot10^{-2},\quad D=4\cdot10^1.
  \]
\end{enumerate}

\subsection*{Сложная задача}

\begin{enumerate}[resume]
  \item Высота полёта камня задаётся формулой
  \[
  y=-\frac{1}{80}x^2+\frac54x.
  \]
  Камень должен перелететь стену высотой \(12\) м с запасом \(3\) м. Найдите наибольшее расстояние от стены, на котором может находиться человек, бросающий камень, если высота камня у стены должна быть равна \(15\) м.
\end{enumerate}

\newpage

\section*{Ответы к тренировочному блоку}

\begin{center}
\renewcommand{\arraystretch}{1.45}
\begin{tabularx}{\linewidth}{@{}cX@{}}
\toprule
\textbf{№} & \textbf{Ответ} \\
\midrule
1 & \(20\) \\
2 & \(3000\) \\
3 & \(1{,}4\) \\
4 & \(15\) \\
5 & \(3000\) \\
6 & \(10\) \\
7 & \(40\) \\
8 & \(3{,}15\) \\
9 & \(5000\) \\
10 & \(60\) \\
\bottomrule
\end{tabularx}
\end{center}

\end{document}